\section{Justificación}
\label{sec:justificacion}

La elaboración de una tesis de posgrado requiere una comunicación constante entre el estudiante y su asesor. Sin embargo, en la práctica diaria este proceso suele llevarse a cabo por medios informales como aplicaciones de mensajería instantánea y correos electrónicos personales. Aunque estas herramientas son de uso común, no fueron diseñadas para la gestión académica. Como resultado, los borradores se mezclan entre tantos mensajes, se borran archivos por error, no se sabe con certeza cuál es la versión más reciente del trabajo y se olvidan las tareas acordadas porque no existe un registro formal de las reuniones.

A su vez, los sistemas informáticos con los que cuenta la institución se enfocan casi por completo en trámites escolares generales (como inscripciones y calificaciones), dejando sin apoyo el seguimiento continuo del trabajo de investigación. Por esta razón, se justifica el diseño y desarrollo de una \textbf{página web} que concentre en un solo lugar digital todo el proceso de asesorías, resolviendo de manera directa las necesidades de cada participante:

\begin{enumerate}
    \item \textbf{Beneficio para los estudiantes:} Les proporciona un espacio seguro donde subir sus avances de tesis en formato PDF, con la tranquilidad de que sus entregas anteriores quedan guardadas en un historial y no se van a perder ni sobreescribir. Además, les permite revisar en cualquier momento las minutas con los compromisos y fechas límite acordados con su profesor.
    
    \item \textbf{Beneficio para los asesores:} Les facilita revisar los documentos ordenadamente por versiones, retroalimentar el trabajo y registrar en pocos minutos los acuerdos al término de cada sesión, evitando confusiones y malos entendidos sobre las tareas pendientes.
    
    \item \textbf{Beneficio para la coordinación de posgrado:} Le permite saber el estado real del avance de los estudiantes sin tener que pedir reportes manuales a cada maestro. Gracias al control de fechas, la coordinación puede identificar de inmediato qué alumnos tienen más de 30 días de inactividad, pudiendo intervenir a tiempo para evitar que abandonen la tesis o se atrasen en su titulación.
\end{enumerate}

Por último, delimitar el alcance a una primera versión centrada en tres funciones básicas (control de archivos PDF sin sobreescritura, minutas digitales y alertas por inactividad mayor a 30 días) garantiza la viabilidad del proyecto. Esto permite concluir el sistema dentro de los tiempos estipulados para la tesina, entregando una herramienta funcional, fácil de utilizar y útil para la facultad, sobre la cual se podrán agregar más funciones en desarrollos futuros.